\documentclass[11pt,a4paper]{article}

\usepackage[T1]{fontenc}
\usepackage[utf8]{inputenc}
\usepackage[main=italian,provide=*]{babel}
\usepackage{amsmath,amssymb}
\usepackage[a4paper,margin=3cm]{geometry}
\usepackage[colorlinks=true,linkcolor=black,citecolor=black,urlcolor=blue!55!black]{hyperref}
\usepackage{setspace}

\pagestyle{empty}

\begin{document}

\begin{center}
{\large\bfseries Università degli Studi di Parma}\\[2pt]
{\normalsize Dipartimento di Scienze Matematiche, Fisiche e Informatiche}\\[4pt]
{\normalsize Corso di Laurea in Fisica}\\[18pt]

{\Large\bfseries Quantum simulation of molecular spin systems\\ on quantum hardware}\\[6pt]
{\normalsize\itshape Simulazione quantistica di sistemi di spin molecolari su hardware quantistico}\\[18pt]

{\normalsize Candidato: \textbf{Samuele Schiavi}}\\[4pt]
{\normalsize Relatore: \textbf{Prof. Alessandro Chiesa}}\\[18pt]
{\normalsize\bfseries Presentazione}
\end{center}

\vspace{6pt}
\onehalfspacing
\noindent
I magneti molecolari costituiscono piattaforme privilegiate per lo
studio di fenomeni quantistici, ma la loro simulazione tramite computer
classici diventa rapidamente intrattabile al crescere del numero di spin
coinvolti. I computer quantistici, mappando ciascuno spin $1/2$ direttamente
su un qubit, offrono un percorso naturale per aggirare questo limite.
Questo lavoro presenta la simulazione quantistica, su simulatori sia ideali
sia affetti da rumore realistico, di sistemi di spin $1/2$ accoppiati
tramite interazione di Heisenberg e di scambio antisimmetrico di
Dzyaloshinskii--Moriya, di sistemi modello di due o tre spin interagenti con
topologia a catena e ad anello. Lo stato fondamentale è determinato con
l'algoritmo variazionale Variational Quantum Eigensolver (VQE), confrontando
un ansatz generico hardware-efficient con ansatz costruiti tenendo conto
delle simmetrie fisiche del problema; l'evoluzione temporale è simulata
tramite decomposizione di Suzuki--Trotter, e le correlazioni dinamiche
spin-spin sono estratte con un circuito di tipo Hadamard test, che combina
la simulazione dell'evoluzione temporale dei qubit del sistema target con un
qubit ausiliario per estrarre le correlazioni. L'intera pipeline è poi
ripetuta sotto un modello di rumore calibrato su hardware quantistico reale,
individuando il numero di passi di Trotter che bilancia l'errore di
approssimazione con quello accumulato dal rumore, e verificandone la
stabilità rispetto a ottimizzazione sotto rumore e a errore di lettura
asimmetrico. I risultati mostrano che un ansatz allineato alle simmetrie
dell'Hamiltoniana riduce drasticamente il numero di parametri necessari
rispetto a un ansatz generico, un vantaggio che permane, con un costo
maggiore, passando da due a tre spin --- un criterio di progettazione
candidato a estendersi verso sistemi molecolari di dimensione superiore.

\vspace{12pt}
\noindent\textbf{Parole chiave:} simulazione quantistica; sistemi di spin;
Variational Quantum Eigensolver (VQE); decomposizione di Suzuki--Trotter;
rumore quantistico; hardware IBM Quantum.

\end{document}
